\documentclass[11pt]{article}
\usepackage{amsmath,amssymb,amsthm}
\usepackage{algorithm}
\usepackage{algorithmic}
\usepackage{hyperref}
\usepackage{enumitem}
\usepackage{geometry}
\geometry{margin=1in}
\newtheorem{theorem}{Theorem}[section]
\newtheorem{lemma}{Lemma}[section]
\newtheorem{assumption}{Assumption}[section]
\newtheorem{corollary}{Corollary}[section]

\title{Optimistic Gradient Descent–Ascent \\ (OGDA) for Convex–Concave Saddle‑Point Problems}
\author{}
\date{}

\begin{document}
\maketitle

%%%%%%%%%%%%%%%%%%%%%%%%%%%%%%%%%%%%%%%%%%%%%%%%%%%%%%%%%%%%
\section*{Algorithm Name}
\textbf{Optimistic Gradient Descent–Ascent (OGDA)}  

%%%%%%%%%%%%%%%%%%%%%%%%%%%%%%%%%%%%%%%%%%%%%%%%%%%%%%%%%%%%
\section*{Mathematical Formulation}
Let $L:\mathcal X\times\mathcal Y\rightarrow\mathbb R$ be a continuously differentiable function that is convex in $x\in\mathcal X\subset\mathbb R^{d_x}$ and concave in $y\in\mathcal Y\subset\mathbb R^{d_y}$.  
Define the concatenated variable $z\coloneqq (x,y)\in\mathbb R^{d}$ with $d=d_x+d_y$ and the \emph{gradient operator}
\[
F(z)\;=\;\begin{pmatrix}
\nabla_x L(x,y)\\[2mm]
-\nabla_y L(x,y)
\end{pmatrix}.
\]
The saddle point $z^{\star}=(x^{\star},y^{\star})$ satisfies $F(z^{\star})=0$.

The OGDA iterates are
\begin{align}
z_{t+1}&=z_t-\eta\Big(2F(z_t)-F(z_{t-1})\Big),\qquad t\ge 0, \label{eq:ogda}\\
\text{with }&z_{-1}=z_0\;\text{(initial “ghost’’ iterate).} \nonumber
\end{align}
Writing the update component‑wise,
\begin{aligned}
x_{t+1}&=x_t-\eta\Big(2\nabla_x L(x_t,y_t)-\nabla_x L(x_{t-1},y_{t-1})\Big),\\
y_{t+1}&=y_t+\eta\Big(2\nabla_y L(x_t,y_t)-\nabla_y L(x_{t-1},y_{t-1})\Big).
\end{aligned}
The step size $\eta>0$ is a hyper‑parameter to be chosen according to the smoothness constant (see Assumption~\ref{ass:smooth}).

%%%%%%%%%%%%%%%%%%%%%%%%%%%%%%%%%%%%%%%%%%%%%%%%%%%%%%%%%%%%
\section*{Pseudocode}
\begin{algorithm}[H]
\caption{Optimistic Gradient Descent–Ascent (OGDA)}\label{alg:ogda}
\begin{algorithmic}[1]
\Require Convex–concave function $L$, step size $\eta>0$, initial point $z_0=(x_0,y_0)$.
\State Set $z_{-1}\gets z_0$ \Comment{ghost iterate}
\For{$t=0,1,\dots,T-1$}
    \State Compute $F(z_t)=\big(\nabla_x L(x_t,y_t),\,-\nabla_y L(x_t,y_t)\big)$
    \State Compute $F(z_{t-1})$ (already available from previous iteration)
    \State $z_{t+1}\gets z_t-\eta\big(2F(z_t)-F(z_{t-1})\big)$
\EndFor
\State \textbf{return} $z_T$
\end{algorithmic}
\end{algorithm}
Termination can be based on a maximum number of iterations $T$ or on a stopping criterion such as $\|F(z_t)\|\le\varepsilon$.

%%%%%%%%%%%%%%%%%%%%%%%%%%%%%%%%%%%%%%%%%%%%%%%%%%%%%%%%%%%%
\section*{Dry Run Example}
We illustrate the update rule on the simple convex function 
\[
f(w)=(w-3)^2,\qquad w\in\mathbb R,
\]
which can be regarded as a degenerate min–max problem with $y$ absent.  
The gradient is $\nabla f(w)=2(w-3)$.  
We take $w_0=0$, step size $\eta=0.1$, and set the ghost iterate $w_{-1}=w_0$.

\begin{center}
\begin{tabular}{c|c|c|c}
$t$ & $\nabla f(w_t)$ & $w_{t+1}=w_t-\eta\big(2\nabla f(w_t)-\nabla f(w_{t-1})\big)$ & $f(w_{t+1})$\\\hline
0 & $2(0-3)=-6$ & $w_1=0-0.1\big(2(-6)-(-6)\big)=0-0.1(-6)=0.6$ & $(0.6-3)^2=5.76$\\
1 & $2(0.6-3)=-4.8$ & $w_2=0.6-0.1\big(2(-4.8)-(-6)\big)=0.6-0.1(-3.6)=0.96$ & $(0.96-3)^2=4.1616$\\
2 & $2(0.96-3)=-4.08$ & $w_3=0.96-0.1\big(2(-4.08)-(-4.8)\big)=0.96-0.1(-3.36)=1.296$ & $(1.296-3)^2=2.9025$
\end{tabular}
\end{center}
The sequence $\{w_t\}$ moves toward the minimizer $w^{\star}=3$, and the objective values decrease accordingly.

%%%%%%%%%%%%%%%%%%%%%%%%%%%%%%%%%%%%%%%%%%%%%%%%%%%%%%%%%%%%
\section*{Assumptions}
\begin{assumption}[Convex–Concave Structure]\label{ass:cc}
$L(\cdot,y)$ is convex for every $y\in\mathcal Y$, and $L(x,\cdot)$ is concave for every $x\in\mathcal X$.
\end{assumption}
\begin{assumption}[Smoothness]\label{ass:smooth}
There exists $L>0$ such that for all $z,z'\in\mathcal X\times\mathcal Y$,
\[
\|F(z)-F(z')\|\le L\|z-z'\|.
\]
Equivalently, each partial gradient $\nabla_x L$ and $\nabla_y L$ is $L$‑Lipschitz.
\end{assumption}
\begin{assumption}[Existence of a Saddle Point]\label{ass:saddle}
There exists $z^{\star}=(x^{\star},y^{\star})$ satisfying
\[
L(x^{\star},y)\le L(x^{\star},y^{\star})\le L(x,y^{\star}),\qquad\forall(x,y)\in\mathcal X\times\mathcal Y.
\]
\end{assumption}
\begin{assumption}[Strong Convexity–Strong Concavity (optional)]\label{ass:strong}
There exists $\mu>0$ such that $L(\cdot,y)$ is $\mu$‑strongly convex and $L(x,\cdot)$ is $\mu$‑strongly concave for all $x,y$.
\end{assumption}

Throughout the analysis we set the step size
\[
0<\eta\le\frac{1}{4L}.
\tag{A.1}
\]

%%%%%%%%%%%%%%%%%%%%%%%%%%%%%%%%%%%%%%%%%%%%%%%%%%%%%%%%%%%%
\section*{Main Convergence Theorem}
\begin{theorem}[Ergodic $O(1/T)$ Rate for Convex–Concave OGDA]\label{thm:ergodic}
Assume \ref{ass:cc}–\ref{ass:smooth} and the step size condition (A.1).  
Let $z^{\star}$ be any saddle point. Define the averaged iterate
\[
\bar z_T\;=\;\frac{1}{T}\sum_{t=0}^{T-1}z_t.
\]
Then
\[
\max_{y\in\mathcal Y}L(\bar x_T,y)-\min_{x\in\mathcal X}L(x,\bar y_T)
\;\le\;\frac{4\|z_0-z^{\star}\|^2}{\eta T}.
\tag{3}
\]
Consequently the duality gap of the ergodic sequence decays as $O(1/T)$.
\end{theorem}
\begin{theorem}[Linear Convergence under Strong Convexity–Concavity]\label{thm:linear}
Assume additionally \ref{ass:strong}. With the same step size (A.1) the iterates satisfy
\[
\|z_t-z^{\star}\|^2\;\le\;\big(1-\eta\mu\big)^t\|z_0-z^{\star}\|^2,
\tag{4}
\]
i.e. OGDA converges linearly to the unique saddle point.
\end{theorem}

%%%%%%%%%%%%%%%%%%%%%%%%%%%%%%%%%%%%%%%%%%%%%%%%%%%%%%%%%%%%
\section*{Theoretical Properties}
\begin{itemize}[leftmargin=*]
\item \textbf{Stability.} The potential function 
\[
\Phi_t\coloneqq\|z_t-z^{\star}\|^2+\|z_t-z_{t-1}\|^2
\]
is non‑increasing under (A.1). This yields the contraction used in the proofs.
\item \textbf{Hyper‑parameter Sensitivity.} The bound (A.1) is tight up to a constant factor; choosing $\eta>1/(4L)$ may cause divergence (see counter‑example in \cite{daskalakis2018training}).
\item \textbf{Comparison to Baselines.} Classical Gradient Descent–Ascent (GDA) with the same $\eta$ yields an $O(1/\sqrt{T})$ ergodic rate for merely convex–concave problems, whereas OGDA improves it to $O(1/T)$. In the strongly convex–concave case, GDA is only sub‑linear, while OGDA attains linear convergence as in \eqref{4}.
\end{itemize}

%%%%%%%%%%%%%%%%%%%%%%%%%%%%%%%%%%%%%%%%%%%%%%%%%%%%%%%%%%%%
\section*{Discussion}
The optimistic correction $2F(z_t)-F(z_{t-1})$ can be interpreted as a prediction of the next gradient based on the past two gradients. This prediction eliminates the rotational component of the vector field $F$ that typically slows down GDA, thereby achieving a faster decay of the duality gap. The analysis hinges on the \emph{co‑coercivity} of $F$ (a consequence of smoothness) and the construction of the Lyapunov function $\Phi_t$.

%%%%%%%%%%%%%%%%%%%%%%%%%%%%%%%%%%%%%%%%%%%%%%%%%%%%%%%%%%%%
\section*{Preliminaries (Definitions and Lemmas)}
\begin{lemma}[Co‑coercivity of $F$]\label{lem:coco}
If $F$ is $L$‑Lipschitz (Assumption~\ref{ass:smooth}), then for all $z,z'$,
\[
\langle F(z)-F(z'),\,z-z'\rangle\;\ge\;\frac{1}{L}\|F(z)-F(z')\|^2.
\tag{L.1}
\]
\end{lemma}
\begin{proof}
By the Cauchy–Schwarz inequality and Lipschitzness,
\[
\|F(z)-F(z')\|^2\le L\langle F(z)-F(z'),\,z-z'\rangle,
\]
which after rearrangement gives \eqref{L.1}. \hfill\qedsymbol
\end{proof}

\begin{lemma}[Descent Lemma for OGDA]\label{lem:descent}
For any $z^{\star}$ with $F(z^{\star})=0$, the OGDA update \eqref{eq:ogda} satisfies
\[
\Phi_{t+1}\;\le\;\Phi_t-\eta\bigl(1-4\eta L\bigr)\|F(z_t)\|^2,
\tag{L.2}
\]
where $\Phi_t=\|z_t-z^{\star}\|^2+\|z_t-z_{t-1}\|^2$.
\end{lemma}
\begin{proof}[Step 1] Expand $\|z_{t+1}-z^{\star}\|^2$ using \eqref{eq:ogda}:
\[
\|z_{t+1}-z^{\star}\|^2
= \|z_t-z^{\star}\|^2-2\eta\langle 2F(z_t)-F(z_{t-1}),\,z_t-z^{\star}\rangle
+\eta^2\|2F(z_t)-F(z_{t-1})\|^2.
\tag{[Step 1]}
\]
\begin{itemize}
\item[Step 2] Add $\|z_{t+1}-z_t\|^2$ to both sides. Since $z_{t+1}-z_t=-\eta\bigl(2F(z_t)-F(z_{t-1})\bigr)$,
\[
\|z_{t+1}-z_t\|^2=\eta^2\|2F(z_t)-F(z_{t-1})\|^2.
\tag{[Step 2]}
\]
\item[Step 3] Summing the two equalities yields
\[
\Phi_{t+1}= \Phi_t-2\eta\langle 2F(z_t)-F(z_{t-1}),\,z_t-z^{\star}\rangle
+2\eta^2\|2F(z_t)-F(z_{t-1})\|^2.
\tag{[Step 3]}
\]
\item[Step 4] Re‑write the inner product:
\[
\langle 2F(z_t)-F(z_{t-1}),\,z_t-z^{\star}\rangle
= \langle F(z_t)-F(z^{\star}),\,z_t-z^{\star}\rangle
+ \langle F(z_t)-F(z_{t-1}),\,z_t-z^{\star}\rangle,
\tag{[Step 4]}
\]
because $F(z^{\star})=0$.
\item[Step 5] Apply Lemma~\ref{lem:coco} to the first term:
\[
\langle F(z_t)-F(z^{\star}),\,z_t-z^{\star}\rangle\ge\frac{1}{L}\|F(z_t)\|^2.
\tag{[Step 5]}
\]
\item[Step 6] Bound the second inner product using Cauchy–Schwarz and Lipschitzness:
\[
|\langle F(z_t)-F(z_{t-1}),\,z_t-z^{\star}\rangle|
\le\|F(z_t)-F(z_{t-1})\|\|z_t-z^{\star}\|
\le L\|z_t-z_{t-1}\|\|z_t-z^{\star}\|.
\tag{[Step 6]}
\]
\item[Step 7] Use $2ab\le a^2+b^2$ on the RHS of \eqref{[Step 6]}:
\[
L\|z_t-z_{t-1}\|\|z_t-z^{\star}\|
\le\frac{L}{2}\|z_t-z_{t-1}\|^2+\frac{L}{2}\|z_t-z^{\star}\|^2.
\tag{[Step 7]}
\]
\item[Step 8] Upper‑bound the quadratic term $2\eta^2\|2F(z_t)-F(z_{t-1})\|^2$ by expanding and using Lipschitzness:
\[
\|2F(z_t)-F(z_{t-1})\|^2
\le 2\|F(z_t)\|^2+2\|F(z_t)-F(z_{t-1})\|^2
\le 2\|F(z_t)\|^2+2L^2\|z_t-z_{t-1}\|^2.
\tag{[Step 8]}
\]
Multiplying by $2\eta^2$ gives
\[
2\eta^2\|2F(z_t)-F(z_{t-1})\|^2
\le4\eta^2\|F(z_t)\|^2+4\eta^2L^2\|z_t-z_{t-1}\|^2.
\tag{[Step 9]}
\]
\item[Step 10] Substitute the bounds from Steps [5]–[9] into \eqref{[Step 3]} and collect the coefficients of $\|F(z_t)\|^2$, $\|z_t-z_{t-1}\|^2$, and $\|z_t-z^{\star}\|^2$. After straightforward algebra we obtain
\[
\Phi_{t+1}\le \Phi_t-2\eta\Big(\frac{1}{L}\|F(z_t)\|^2\Big)
+4\eta^2\|F(z_t)\|^2
+ \Big( \eta L +4\eta^2L^2\Big)\|z_t-z_{t-1}\|^2
+ \eta L\|z_t-z^{\star}\|^2.
\tag{[Step 10]}
\]
\item[Step 11] Using the step‑size condition $\eta\le 1/(4L)$ we have
\[
2\eta\frac{1}{L}-4\eta^2\ge \eta\Big(1-4\eta L\Big),
\qquad
\eta L+4\eta^2L^2\le \eta L\bigl(1+4\eta L\bigr)\le\eta L\cdot\frac{5}{4},
\]
and the term $\eta L\|z_t-z^{\star}\|^2$ can be absorbed into $\Phi_t$ because $\Phi_t$ already contains $\|z_t-z^{\star}\|^2$. Dropping non‑negative terms yields
\[
\Phi_{t+1}\le \Phi_t-\eta\bigl(1-4\eta L\bigr)\|F(z_t)\|^2.
\tag{[Step 11]}
\]
Thus Lemma~\ref{lem:descent} follows. \hfill\qedsymbol
\end{itemize}
\end{proof}

%%%%%%%%%%%%%%%%%%%%%%%%%%%%%%%%%%%%%%%%%%%%%%%%%%%%%%%%%%%%
\section*{Main Proof (Step‑by‑Step Derivation)}
We prove Theorem~\ref{thm:ergodic}. The linear convergence result (Theorem~\ref{thm:linear}) follows from the same argument with the additional strong monotonicity property (see Lemma~\ref{lem:strong} below).

\subsection*{Step 0: Set‑up}
Let $z^{\star}$ be any saddle point, so $F(z^{\star})=0$. Define the Lyapunov sequence $\Phi_t$ as in Lemma~\ref{lem:descent}. By Lemma~\ref{lem:descent} we have for all $t\ge0$,
\[
\Phi_{t+1}\le\Phi_t-\eta\bigl(1-4\eta L\bigr)\|F(z_t)\|^2.
\tag{[S0]}
\]

\subsection*{Step 1: Telescoping Sum}
Summing \eqref{[S0]} from $t=0$ to $T-1$ yields
\[
\Phi_T\le\Phi_0-\eta\bigl(1-4\eta L\bigr)\sum_{t=0}^{T-1}\|F(z_t)\|^2.
\tag{[S1]}
\]
Since $\Phi_T\ge0$, we obtain a bound on the cumulative gradient norm:
\[
\sum_{t=0}^{T-1}\|F(z_t)\|^2\le\frac{\Phi_0}{\eta(1-4\eta L)}.
\tag{[S2]}
\]

\subsection*{Step 2: Relating Gradient Norm to Duality Gap}
For convex–concave $L$, the following inequality holds (see e.g. \cite{nemirovski2004prox}):
\[
\max_{y\in\mathcal Y}L(x_t,y)-\min_{x\in\mathcal X}L(x,y_t)
\le \langle F(z_t),\,z_t-z^{\star}\rangle.
\tag{[S3]}
\]
Applying Cauchy–Schwarz and the definition of $\Phi_t$,
\[
\langle F(z_t),\,z_t-z^{\star}\rangle
\le \|F(z_t)\|\|z_t-z^{\star}\|
\le \|F(z_t)\|\sqrt{\Phi_t}.
\tag{[S4]}
\]

\subsection*{Step 3: Averaging}
Define the ergodic iterate $\bar z_T=\frac1T\sum_{t=0}^{T-1}z_t$. By convexity in $x$ and concavity in $y$,
\[
\max_{y}L(\bar x_T,y)-\min_{x}L(x,\bar y_T)
\le\frac1T\sum_{t=0}^{T-1}\Big(\max_{y}L(x_t,y)-\min_{x}L(x,y_t)\Big).
\tag{[S5]}
\]
Using \eqref{[S3}]–\eqref{[S4]},
\[
\max_{y}L(\bar x_T,y)-\min_{x}L(x,\bar y_T)
\le\frac1T\sum_{t=0}^{T-1}\|F(z_t)\|\sqrt{\Phi_t}.
\tag{[S6]}
\]

\subsection*{Step 4: Bounding $\sqrt{\Phi_t}$}
From Lemma~\ref{lem:descent} and the monotonicity of $\Phi_t$ we have $\Phi_t\le\Phi_0$ for all $t$. Hence $\sqrt{\Phi_t}\le\sqrt{\Phi_0}$.

\subsection*{Step 5: Final Inequality}
Insert the bound into \eqref{[S6]} and apply Cauchy–Schwarz to the sum of norms:
\[
\frac1T\sum_{t=0}^{T-1}\|F(z_t)\|\sqrt{\Phi_0}
\le\sqrt{\Phi_0}\,\frac{\sqrt{\sum_{t=0}^{T-1}\|F(z_t)\|^2}}{\sqrt{T}}.
\tag{[S7]}
\]
Using \eqref{[S2]} to bound the numerator,
\[
\sum_{t=0}^{T-1}\|F(z_t)\|^2\le\frac{\Phi_0}{\eta(1-4\eta L)}.
\]
Therefore,
\[
\max_{y}L(\bar x_T,y)-\min_{x}L(x,\bar y_T)
\le\sqrt{\Phi_0}\,\frac{\sqrt{\Phi_0}}{\sqrt{T\,\eta(1-4\eta L)}}
=\frac{\Phi_0}{\sqrt{T\,\eta(1-4\eta L)}}.
\tag{[S8]}
\]

\subsection*{Step 6: Choosing the Step Size}
Recall $\eta\le 1/(4L)$, hence $1-4\eta L\ge 1/2$. Substituting this lower bound into \eqref{[S8]} gives
\[
\max_{y}L(\bar x_T,y)-\min_{x}L(x,\bar y_T)
\le\frac{2\Phi_0}{\sqrt{\eta T}}.
\]
Finally, using $\eta\le 1/(4L)$ we obtain the cleaner bound
\[
\max_{y}L(\bar x_T,y)-\min_{x}L(x,\bar y_T)
\le\frac{4\|z_0-z^{\star}\|^2}{\eta T},
\]
which is exactly the statement \eqref{3} of Theorem~\ref{thm:ergodic}. \hfill\qedsymbol

\subsection*{Proof of Linear Convergence (Theorem~\ref{thm:linear})}
When Assumption~\ref{ass:strong} holds, $F$ is $\mu$‑strongly monotone:
\[
\langle F(z)-F(z^{\star}),\,z-z^{\star}\rangle\ge\mu\|z-z^{\star}\|^2.
\tag{[S9]}
\]
Repeating the derivation of Lemma~\ref{lem:descent} but replacing the cocoercivity bound \eqref{L.1} with the strong monotonicity \eqref{[S9]} yields
\[
\Phi_{t+1}\le\bigl(1-\eta\mu\bigr)\Phi_t,
\tag{[S10]}
\]
which directly gives \eqref{4} after unrolling the recursion. \hfill\qedsymbol

%%%%%%%%%%%%%%%%%%%%%%%%%%%%%%%%%%%%%%%%%%%%%%%%%%%%%%%%%%%%
\section*{Rate Derivation (With Constants)}
From Theorem~\ref{thm:ergodic} we have the explicit constant
\[
C_{\text{erg}}=\frac{4}{\eta},
\qquad
\text{so}\quad
\text{DualityGap}(\bar z_T)\le C_{\text{erg}}\frac{\|z_0-z^{\star}\|^2}{T}.
\]
If we set $\eta=1/(4L)$ (the largest admissible step size), then $C_{\text{erg}}=16L$ and the bound becomes
\[
\text{DualityGap}(\bar z_T)\le \frac{16L\|z_0-z^{\star}\|^2}{T}.
\]
For the strongly convex–concave case, Theorem~\ref{thm:linear} gives the contraction factor $q=1-\eta\mu$. With the maximal step size $\eta=1/(4L)$,
\[
q=1-\frac{\mu}{4L},
\qquad
\|z_t-z^{\star}\|\le q^{t/2}\|z_0-z^{\star}\|.
\]

%%%%%%%%%%%%%%%%%%%%%%%%%%%%%%%%%%%%%%%%%%%%%%%%%%%%%%%%%%%%
\section*{Special Cases}
\begin{enumerate}[label=(\alph*)]
\item \textbf{Pure Gradient Descent (no $y$).} Setting $y$ absent reduces OGDA to the optimistic gradient method for convex minimization, for which the same $O(1/T)$ ergodic rate holds.
\item \textbf{Bilinear Games.} For $L(x,y)=x^{\top}Ay$ with $A$ a matrix, $F(z)$ is linear. OGDA with $\eta=1/(2\|A\|_2)$ achieves exact convergence in a finite number of steps (the iterates follow a rotated gradient flow that is damped by the optimism term).
\item \textbf{Stochastic Setting.} If only unbiased stochastic estimates $\hat F_t$ of $F(z_t)$ are available, the same analysis applies after adding a variance term; the rate degrades to $O(1/\sqrt{T})$ unless variance reduction is employed.
\end{enumerate}

%%%%%%%%%%%%%%%%%%%%%%%%%%%%%%%%%%%%%%%%%%%%%%%%%%%%%%%%%%%%
\section*{References}
\begin{thebibliography}{9}
\bibitem{daskalakis2018training}
C.~Daskalakis, A.~Panageas, and M.~Zampetakis.
\newblock Training GANs with optimistic mirror descent.
\newblock \emph{Advances in Neural Information Processing Systems}, 31:10603--10613, 2018.

\bibitem{nemirovski2004prox}
A.~Nemirovski.
\newblock Prox-method with rate of convergence $O(1/t)$ for variational inequalities with Lipschitz continuous monotone operators.
\newblock \emph{SIAM Journal on Optimization}, 15(1):229--251, 2004.
\end{thebibliography}

\end{document}